\documentclass[11pt]{article}
\usepackage[margin=1in]{geometry}
\usepackage[T1]{fontenc}
\usepackage{lmodern,amsmath,amssymb,amsthm,booktabs,longtable,array,enumitem}
\usepackage[colorlinks=true,linkcolor=blue,citecolor=blue,urlcolor=blue]{hyperref}
\newtheorem{theorem}{Theorem}
\newtheorem{proposition}[theorem]{Proposition}
\newtheorem{corollary}[theorem]{Corollary}
\newtheorem{lemma}[theorem]{Lemma}
\theoremstyle{definition}
\newtheorem{definition}[theorem]{Definition}
\newtheorem{example}[theorem]{Counterexample}
\theoremstyle{remark}
\newtheorem{remark}[theorem]{Remark}
\newcommand{\Prb}{\mathbb P}
\newcommand{\E}{\mathbb E}
\newcommand{\ind}{\mathbf 1}
\newcommand{\TV}{d_{\mathrm{TV}}}
\title{Beyond Constant Fault Rates:\\Reliability Bounds for Gated Computation\\and Semantic Mechanization}
\author{Ivo Matija\v{s}evi\'c}
\date{October 2026}
\newcommand\blfootnote[1]{\begingroup\renewcommand\thefootnote{}\footnote{#1}\addtocounter{footnote}{-1}\endgroup}
\begin{document}
\maketitle
\blfootnote{\textcopyright~2026 Ivo Matija\v{s}evi\'c.\ This work is licensed under a Creative Commons
Attribution 4.0 International (CC BY 4.0) license, \url{https://creativecommons.org/licenses/by/4.0/}.\ Preprint, DOI \href{https://doi.org/10.5281/zenodo.23293740}{10.5281/zenodo.23293740}.}
\begin{abstract}
Constant fault intensity, independent segments, fresh retries, and a fixed fault distribution
are different assumptions. Earlier papers on gated computation use different subsets of them.
We separate these assumptions and examine which conclusions survive their removal.
A conditional survival identity replaces the independent-segment product. Uniform conditional
gate accuracy and progress bounds preserve logarithmic verification overhead without a Poisson
arrival model. A verification floor yields a finite horizon only when its exposure persists
along correct prefixes. Unequal difficulty, latent regimes, persistent retry defects, and
changing input distributions give counterexamples to predictions based on one average rate.
Two positive results locate the $1/\sigma$ ceiling multiplier of the first paper. For the ideal
gated survival curve, a latent rate drawn once per run preserves it exactly, and a monotone hazard
orders it: at most $1/\sigma$ for an increasing failure rate and at least $1/\sigma$ for a
decreasing one. With finite segments, the upper bounds hold after rounding to whole segments.
For semantic mechanization, local escape and fallback probabilities replace one volume-weighted
closure fraction. The twin result survives changes in the arrival law and dependent or adaptive
sampling, provided twins induce the same observation law. We classify all 25 named results of
the four earlier papers as arrival-free, conditionally preserved, or stationary-parameter
formulas that need replacement, and we give executable finite-model checks. These are mathematical models, not
measurements of deployed agents. The contribution is their application and synthesis; the
probability identities and inequalities are classical.
\end{abstract}

\section{Introduction}
Part I models faults by a homogeneous Poisson process and derives a verification threshold,
a stealth-limited horizon, and a checkpoint interval~\cite{p1}. Part II studies representation
blind spots~\cite{p2}. Part III studies accepted-output escape, cost, and shared
provenance~\cite{p3}. Part IV studies progressive semantic mechanization (PSM): moving work
from executor calls to deterministic mechanisms~\cite{p4}.

A natural question is whether their conclusions depend on faults entering at a constant rate.
The answer differs by result. A blind-set inclusion is a statement about what a gate can
distinguish. A horizon is a statement about how often unresolved faults occur over a sequence.
One does not determine the other. Even a constant expected count per unit work need not give
an exponential probability of completing without fault.

Ord's analysis motivates the constant-hazard baseline on a particular task suite and leaves
its wider validity open~\cite{ord}. A fit of success against task duration also does not,
by itself, identify within-run fault arrivals, retry behavior, or checkpoint completeness.
We do not select an empirically correct alternative here. We identify the assumptions needed
to use each result and provide countermodels when those assumptions are absent.

The contributions are: (i) a catalog of candidate fault-introduction models, mapped to the
results they affect; (ii) conditional versions of the threshold and floor arguments;
(iii) a cumulative-exposure extension, two results on when the $1/\sigma$ ceiling multiplier
survives, and counterexamples to average-rate substitution;
(iv) a separation of PSM's closure, evidence, drift, and cost assumptions; and
(v) an assessment and classification of every named theorem, proposition, lemma, and corollary
in Parts I--IV, with a list of errata candidates found on the way.
We call a result \emph{preserved} only with its remaining hypotheses. A failed extension does
not refute a theorem inside its original model. Separate wording issues are identified as such.

\section{What is allowed to change?}\label{sec:axes}
We distinguish six axes.
\begin{enumerate}[leftmargin=*,itemsep=2pt]
\item \textbf{Difficulty:} a segment's prior fault probability can depend on its work, length,
and context load. A deterministic rate $\lambda(t)$ is one model.
\item \textbf{Dependence:} segments can share a latent regime, input, mechanism, or upstream fault.
Constant marginal probabilities do not remove this dependence.
\item \textbf{Retry:} a rejected attempt can change the next attempt's distribution. A rollback
can restore an artifact while retaining a faulty premise or environment.
\item \textbf{Fault type and visibility:} the distribution of produced faults and the gates'
blind sets can change together. Fault intensity alone omits this interaction.
\item \textbf{Intervention:} repair, reset, routing, and mechanization can change later state.
Their effect is not generally a substitution into a fixed sequence of risks.
\item \textbf{Semantics and workload:} input distributions, specifications, and mechanisms can
drift. A deterministic mechanism can then have time-varying error without random execution.
\end{enumerate}

Three clocks must also stay separate: useful work or emitted units, attempted work including
retries, and calendar time. The first indexes reliability below. The second determines execution
cost. The third matters for drift and maintenance. A median horizon on one clock is not a
median on another. A fixed semantic defect is not a new arrival on every affected output.

\paragraph{Candidate fault-introduction models.}
Table~\ref{tab:catalog} lists candidate replacements for the constant-rate model. Each is an
established modeling family; none is selected here as empirically correct. The last column names
the earlier results most sensitive to each model; Section~\ref{sec:audit} gives the full mapping.
Adversarial fault introduction is not a stochastic process and is out of scope, as it is for
Part~I's benign-executor analysis.

\begin{table}[h]\centering\small
\caption{Candidate fault-introduction models and where this paper treats them.}\label{tab:catalog}
\begin{tabular}{>{\raggedright\arraybackslash}p{0.265\textwidth}>{\raggedright\arraybackslash}p{0.19\textwidth}>{\raggedright\arraybackslash}p{0.17\textwidth}>{\raggedright\arraybackslash}p{0.27\textwidth}}
\toprule Model & What changes & Treated in & Most sensitive earlier results\\\midrule
Nonhomogeneous Poisson, deterministic $\lambda(t)$, incl.\ Weibull~\cite{goel} & difficulty by position
 & Props.~\ref{prop:hazard}, \ref{prop:monotone} & ceiling multiplier, checkpoint interval, R-A\\
Mixed Poisson / frailty: rate drawn once per run~\cite{vaupel} & dependence via a latent rate
 & Ex.~\ref{ex:latent}, Prop.~\ref{prop:mixing} & horizon value; the multiplier only at finite $s$ (ideal limit: exact)\\
Markov-modulated / regime switching~\cite{mmpp} & dependence, persistence
 & Eq.~\eqref{eq:kernel} & slip, threshold uniformity\\
Self-exciting / propagation~\cite{hawkes} & faults beget faults
 & Sec.~\ref{sec:cost}, propagation & after an escape: damage only; before an escape: slip, horizon, retry cost\\
Common-cause / shared defect & one fault in many outputs
 & Ex.~\ref{ex:shared} & series identity, horizon\\
Persistent retry / imperfect repair~\cite{brownproschan} & retry distribution
 & Prop.~\ref{prop:retry}, Ex.~\ref{ex:retry}, Eq.~\eqref{eq:kernel} & slip, checkpoint cost, R-C\\
Load- or context-dependent hazard & hazard as a function of state
 & Prop.~\ref{prop:monotone}; R-E row & R-E, ceiling multiplier\\
Visibility that co-varies with intensity & fault type $\times$ gate
 & Sec.~\ref{sec:hazard}, Prop.~\ref{prop:mixing} & diversity, ceiling\\
Input or specification drift & the fault set itself
 & Prop.~\ref{prop:drift} & replay, evidence, validation\\
\bottomrule
\end{tabular}
\end{table}

\begin{definition}[Success and emission]\label{def:success}
Fix a workload and policy. Unless stated otherwise, each requested unit eventually emits with
probability one. Let $I_i$ indicate that emitted unit $i$ has an escaped fault. Put
$A_n=\{I_1=\cdots=I_n=0\}$, $S_n=\Prb(A_n)$, and $A_0=\Omega$.
Success means no escaped fault in the first $n$ emitted units. Later repair does not erase an
earlier escape in this metric. The discrete half-success horizon is
$H_{1/2}=\inf\{n:S_n\le1/2\}$, with $\inf\varnothing=\infty$.
\end{definition}
If emission can stall, report its probability and cost separately, or count non-completion as
failure. Conditioning on eventual completion can select easier runs. We do not use that
conditioning to conceal failures. Damage, the number of faulty outputs, and terminal correctness
after repair are different metrics. They require additional models.

\section{Conditional risk replaces constant rate}
\begin{proposition}[Survival and exposure]\label{prop:survival}
For each $i$ with $S_{i-1}>0$, define $q_i=\Prb(I_i=1\mid A_{i-1})$. Then
\begin{equation}\label{eq:survival}
S_n=\prod_{i=1}^n(1-q_i),\qquad
-\log S_n=\sum_{i=1}^n-\log(1-q_i).
\end{equation}
The second equality uses $-\log0=\infty$. After survival reaches zero, set later $q_i=0$.
If $\ell_i\le q_i\le u_i$ at every index where $q_i$ is defined, with $0\le\ell_i\le u_i\le1$, then
\[
\prod_i(1-u_i)\le S_n\le\prod_i(1-\ell_i),\qquad
1-S_n\le\sum_i q_i\le\sum_i u_i.
\]
If all $q_i<1$, then $S_\infty>0$ if and only if $\sum_i q_i<\infty$.
\end{proposition}
\begin{proof}
The chain rule gives $S_i=S_{i-1}(1-q_i)$. Iterate and take logarithms.
The product bounds follow factor by factor. The inequality
$1-\prod_i(1-q_i)\le\sum_iq_i$ follows by induction.
For the last claim, $-\log(1-x)\ge x$. Conversely, summability makes $q_i\to0$;
eventually $q_i\le1/2$, where $-\log(1-q_i)\le2q_i$. The finite prefix has positive product.
\end{proof}

The identity is the chain rule, not a new independence theorem. Its useful consequence is
where bounds must hold: on histories that have remained correct. Unconditional marginals
$f_i=\Prb(I_i=1)$ cannot generally replace $q_i$.
For arbitrary dependence they imply only
\begin{equation}\label{eq:marginal}
\max_i f_i\le\Prb\Big(\bigcup_i\{I_i=1\}\Big)\le\min(1,\sum_i f_i).
\end{equation}
These bounds are sharp: nested events attain the lower bound; disjoint events attain the upper
bound when the sum is at most one. No homogeneous arrival process is needed for either bound.

\begin{example}[Same marginal risk, different horizons]\label{ex:shared}
Let $f=0.01$. Independent indicators give $S_{100}=0.99^{100}\approx0.3660$.
If $I_i=Z$ for all $i$, with $Z$ Bernoulli$(0.01)$, then $S_n=0.99$ for every $n\ge1$.
Here $q_1=0.01$ and $q_i=0$ for $i>1$. Both models have expected fault count $0.01n$.
The second has no finite half-success horizon, but has probability $0.01$ of all $n$ outputs
being faulty. Dependence changes the distribution of damage, not just a horizon constant.
\end{example}

\begin{example}[A floor that declines with exposure]\label{ex:summable}
Let independent segment slips have $q_i=0.1\,2^{-i}$. Every segment has positive risk, yet
$\Prb(\text{any escape ever})\le\sum_{i\ge1}q_i=0.1$. Thus a positive local floor at each
step does not imply an eventual escape with probability one, or a finite median horizon.
A persistent lower bound, or sufficient cumulative exposure, is needed.
\end{example}

\section{Retries and a threshold without Poisson arrivals}
\begin{proposition}[First accepted attempt]\label{prop:retry}
For a single unit, let $R_r$ mean attempt $r$ is reached, with $R_1=\Omega$.
Let $a_r=\Prb(\text{accept at }r\mid R_r)$ and
$b_r=\Prb(\text{bad and accept at }r\mid R_r)$. If emission occurs almost surely, then
\begin{equation}\label{eq:retry}
\Prb(\text{emitted bad})=\sum_{r\ge1}\Prb(R_r)b_r,
\quad \E N=\sum_{r\ge1}\Prb(R_r),
\quad \Prb(R_{r+1})=\Prb(R_r)(1-a_r).
\end{equation}
If $a_r=a>0$ and $b_r=b$ for every reached attempt, these reduce to $b/a$ and $1/a$.
More generally, if $a_r\ge a_0>0$ and $b_r\le b_0$, then
$\E N\le1/a_0$ and $\Prb(\text{emitted bad})\le b_0/a_0$.
A ratio bound is often sharper: if emission occurs almost surely and $b_r\le\kappa a_r$ at
every reached attempt, then $\Prb(\text{emitted bad})\le\min(1,\kappa)$.
\end{proposition}
\begin{proof}
First acceptance events are disjoint. Sum their bad parts. The tail-sum identity gives $\E N$.
Conditional rejection gives the recursion. A geometric tail follows from $a_r\ge a_0$.
Summing that tail proves both bounds and also establishes almost-sure emission in this case.
For the ratio bound, $\sum_r\Prb(R_r)b_r\le\kappa\sum_r\Prb(R_r)a_r=\kappa$, because the last
sum is the emission probability.
\end{proof}
The same reasoning applies conditional on a pre-unit history. Independent and identically
distributed retries are sufficient for constant $a,b$, but are not necessary.
In the usual notation, $a=(1-p)(1-\beta_m)+p\alpha_m$ and $b=p\alpha_m$.

\begin{example}[Persistent difficulty invalidates pooled retry arithmetic]\label{ex:retry}
Choose a unit type once. With probability $1/2$, it is always good and accepted.
Otherwise it is always bad and each retry accepts independently with probability $1/2$.
Every unit eventually emits. A fresh, pooled first attempt has $a=3/4$, $b=1/4$, but
the eventual bad fraction is $1/2$, not $b/a=1/3$.
Expected attempts are $\tfrac12(1)+\tfrac12(2)=3/2$, not $1/a=4/3$.
If bad units instead always reject, half the units never emit. A fresh first-attempt
acceptance probability of $1/2$ does not imply finite expected completion time.
\end{example}

\begin{theorem}[Uniform conditional threshold]\label{thm:threshold}
Fix $n$ units and $0<\delta<1$. At every reachable pre-attempt history, including previous
rejections and earlier slips, suppose:
\begin{enumerate}[label=(\roman*),itemsep=1pt]
\item the candidate is bad with probability at most $\bar p<1$;
\item conditional on its good/bad status and that history, the $m$ gate-error indicators
are independent, each with probability at most $1/2-\gamma$, for fixed $\gamma>0$;
\item a majority decides acceptance; an attempt costs at most $s+mg$ for fixed $s,g$.
\end{enumerate}
Choose odd $m$ and write $d_m=e^{-2m\gamma^2}$. If
\[
d_m\le\min\{1/2,\ \delta(1-\bar p)/(2n\bar p)\},
\]
where the second constraint is omitted when $\bar p=0$, then
\[
S_n\ge1-\delta,\qquad
\E[\text{total cost}]\le\frac{2n(s+mg)}{1-\bar p}.
\]
For fixed $\bar p,\gamma,s,g$ this gives $m=O(\log(n/\delta))$ and logarithmic
multiplicative overhead. Neither Poisson arrivals nor independence across units is required.
\end{theorem}
\begin{proof}
Hoeffding's inequality bounds each majority error by $d_m$~\cite{hoeffding}.
At every history, acceptance has probability at least
$(1-\bar p)(1-d_m)\ge(1-\bar p)/2$, and bad acceptance has probability at most $\bar p d_m$.
These uniform bounds remain valid after averaging over rejected histories.
The ratio bound of Proposition~\ref{prop:retry}, with $\kappa=2\bar p d_m/(1-\bar p)$, bounds
each unit's slip conditional on every pre-unit history by $\kappa\le\delta/n$. Its geometric
tail bounds the expected attempts by $2/(1-\bar p)$.
Average the slip bound over correct prefixes and apply Proposition~\ref{prop:survival}.
Sum conditional expected costs over all units, including histories with earlier slips.
\end{proof}

This theorem permits changing priors and adaptive history. Its uniform hypotheses are strong.
A latent hard state can violate (i); shared gate blindness can violate (ii). Bounds only on
correct prefixes suffice for the reliability calculation, but not for unconditional expected
cost after a slip. A decreasing gate edge also changes the conclusion: the sufficient gate count
scales as $\gamma^{-2}\log(n/\delta)$, not uniformly as $\log n$ if $\gamma$ tends to zero.

\begin{proposition}[Conditional verification floor]\label{prop:floor}
At each correct pre-unit history $h$, suppose the first attempt is bad with probability $p_i(h)$.
Conditional on badness, a fraction $\sigma_i(h)$ passes every gate and the aggregator accepts
unanimous passes. Then
\[
q_i\ge \E[p_i(h)\sigma_i(h)\mid A_{i-1}].
\]
In particular, if $p_i(h)\sigma_i(h)\ge\ell_i$ uniformly, then
$S_n\le\prod_i(1-\ell_i)\le\exp(-\sum_i\ell_i)$.
If $\ell_i\ge\ell>0$, then
$H_{1/2}\le\lceil\log 2/[-\log(1-\ell)]\rceil$ for $\ell<1$.
\end{proposition}
\begin{proof}
A stealth fault on the first attempt immediately emits a bad unit. Subsequent attempts can
only add to this event, not remove it. Average over correct histories and use
Proposition~\ref{prop:survival} and $1-x\le e^{-x}$.
\end{proof}
Fresh retries are unnecessary for this lower bound. Persistent exposure on correct histories
is necessary for the stated horizon consequence. The shared-defect example has no such
exposure after its first correct output.

\section{Unequal difficulty and cumulative hazard}\label{sec:hazard}
For a deterministic locally integrable nonnegative intensity $\lambda(t)$, a nonhomogeneous
Poisson process has no-fault probability $e^{-\Lambda(T)}$, where
$\Lambda(T)=\int_0^T\lambda(t)\,dt$. Its raw half-success horizon is the first time
$\Lambda$ reaches $\log2$, if it ever does. This extends the baseline without adding dependence.
An intensity depending on random history is not interchangeable with this deterministic function.

\begin{proposition}[A vanishing-segment extension of the stealth ceiling]\label{prop:hazard}
Partition $[0,T]$ into segments with cumulative intensities
$L_i=\int_{t_{i-1}}^{t_i}\lambda(t)\,dt$ and raw badness $p_i=1-e^{-L_i}$.
Suppose segments are independent; within each segment, retries are fresh, good candidates
always pass, and a bad candidate passes with probability $\sigma\in(0,1]$.
Then
\begin{equation}\label{eq:finitehazard}
S(T)=\prod_i\frac{1-p_i}{1-p_i+p_i\sigma}.
\end{equation}
As $\max_i L_i\to0$, at fixed finite $\Lambda(T)$,
$S(T)\to e^{-\sigma\Lambda(T)}$. Thus the limiting gated horizon is
$H_\sigma=\inf\{t:\Lambda(t)\ge\log2/\sigma\}$, with $\inf\varnothing=\infty$; when it is finite,
$\Lambda(H_\sigma)=\log2/\sigma$. A bounded $\Lambda$ can give a finite raw horizon and an infinite
gated one: $\lambda(t)=e^{-t}$ and $\sigma=1/2$ give $e^{-\sigma\Lambda(t)}>e^{-1/2}>1/2$. If $\Lambda(t)=(t/\eta)^k$, with $k,\eta>0$, the limiting ratio is
\begin{equation}\label{eq:weibull}
H_{\mathrm{gated}}/H_{\mathrm{raw}}=\sigma^{-1/k}.
\end{equation}
For $\sigma<1$, the ratio $1/\sigma$ is recovered only at $k=1$ in this family.
\end{proposition}
\begin{proof}
Proposition~\ref{prop:retry} gives accepted slip
$p_i\sigma/(1-p_i+p_i\sigma)$. Independent segments give the displayed product.
Its $i$th logarithm is $-L_i-\log(1-(1-\sigma)(1-e^{-L_i}))
=-\sigma L_i+O(L_i^2)$, uniformly as $\max_iL_i\to0$.
Since $\sum_i L_i^2\le(\max_iL_i)\Lambda(T)$, the error vanishes.
Solve the limiting survival curve at one half. For the power cumulative intensity,
$H_{\mathrm{raw}}=\eta(\log2)^{1/k}$ and
$H_{\mathrm{gated}}=\eta(\log2/\sigma)^{1/k}$.
\end{proof}
\begin{corollary}[Finite-segment bound under fresh Poisson retries]\label{cor:finitefloor}
In the same independent-segment model, allow false rejection and any bad-accept probability
$\alpha_i\ge\sigma_i\in[0,1]$. Let acceptance have positive probability. Then
\[
S(T)\le\prod_i\frac{1-p_i}{1-p_i+p_i\sigma_i}
\le\exp\left(-\sum_i\sigma_i L_i\right).
\]
For constant intensity $\lambda$ and stealth share $\sigma>0$, this yields the exact bound
$H\le\log2/(\lambda\sigma)$ at any fixed segment length, under log-linear interpolation between
boundaries: for $T=(k+\theta)s$ with $0\le\theta<1$,
$S(T)=\big[\prod_{i=1}^{k}(1-q_i)\big](1-q_{k+1})^\theta$. With identical gate behavior this is
$(1-q)^{T/s}$, Part I's horizon convention. Other interpolation rules
need not satisfy it: with $\lambda=1$, $\sigma=1/2$, and $s=4$, linear interpolation crosses one
half near $2.0746>2\log2$.
For the discrete horizon, round the number of segments up.
\end{corollary}
\begin{proof}
The accepted slip $p_i\alpha_i/[(1-p_i)(1-\beta_i)+p_i\alpha_i]$ increases with
$\alpha_i$ and $\beta_i$. Set $\alpha_i=\sigma_i$ and $\beta_i=0$ for its lower bound.
The corresponding survival factor is $1/[1-\sigma_i+\sigma_i e^{L_i}]$.
Convexity of the exponential gives $1-\sigma_i+\sigma_i e^{L_i}\ge e^{\sigma_i L_i}$.
Multiply the factors.
\end{proof}
This is a reliability limit, not a cost-free policy: fixed positive gate cost makes arbitrarily
short segments expensive. It models ideal visibility outside the stealth set and zero false
rejection. A time-varying $\sigma(t)$ gives $\exp(-\int_0^T\lambda(t)\sigma(t)dt)$ under
piecewise-continuous bounded marks and the same refinement argument. Independent thinning of
the underlying Poisson process gives that survival exactly, but thinning and the finite-segment
retry protocol are different models.

The relevant average is $\int\lambda\sigma/\int\lambda$, not generally the ordinary
time-average of $\sigma$. On two equal intervals, rates $(0.01,0.09)$ and stealth shares
$(0.1,0.5)$ give mean escape intensity $0.023$. Multiplying their separate means gives $0.015$.
High fault intensity and low visibility can coincide.

The Weibull ratio~\eqref{eq:weibull} is one instance of a general ordering. We use the standard
reliability classes: increasing failure rate (IFR) and decreasing failure rate (DFR)~\cite{bp}.

\begin{proposition}[Monotone hazard orders the multiplier]\label{prop:monotone}
Let $\Lambda$ be continuous and nondecreasing with $\Lambda(0)=0$. For $\sigma\in(0,1]$ put
$H_{\mathrm{raw}}=\inf\{t:\Lambda(t)\ge\log2\}$ and $H_\sigma=\inf\{t:\Lambda(t)\ge\log2/\sigma\}$.
If $\lambda$ is nondecreasing (IFR), then $H_\sigma\le H_{\mathrm{raw}}/\sigma$.
If $\lambda$ is nonincreasing (DFR), then $H_\sigma\ge H_{\mathrm{raw}}/\sigma$.
Under IFR and the model of Corollary~\ref{cor:finitefloor} with $\sigma_i=\sigma$ and equal
segment length $s$, the discrete gated horizon, measured in work, is at most
$s\lceil H_\sigma/s\rceil\le s\lceil H_{\mathrm{raw}}/(\sigma s)\rceil$.
\end{proposition}
\begin{proof}
Under IFR, $\Lambda$ is convex with $\Lambda(0)=0$, so $\Lambda(t)/t$ is nondecreasing for $t>0$.
If $H_{\mathrm{raw}}<\infty$, continuity gives $\Lambda(H_{\mathrm{raw}})=\log2$, and at
$t=H_{\mathrm{raw}}/\sigma\ge H_{\mathrm{raw}}$ we get $\Lambda(t)\ge\Lambda(H_{\mathrm{raw}})/\sigma
=\log2/\sigma$. Under DFR, $\Lambda(t)/t$ is nonincreasing. For
$H_{\mathrm{raw}}\le t<H_{\mathrm{raw}}/\sigma$, $\Lambda(t)\le t\log2/H_{\mathrm{raw}}<\log2/\sigma$;
for $t<H_{\mathrm{raw}}$, $\Lambda(t)<\log2\le\log2/\sigma$. For the last claim,
Corollary~\ref{cor:finitefloor} gives $S(T)\le e^{-\sigma\Lambda(T)}$ at every segment boundary.
At the first boundary $T\ge H_\sigma$, this is at most one half. The rounding is necessary:
with constant rate $1$, $\sigma=1/2$, and $s=1$, one segment survives with probability
$2/(1+e)>1/2$, so the discrete horizon is $2>2\log2$.
\end{proof}
Part~IV's load-monotonicity hypothesis, together with a nondecreasing load (its R-E case (a)),
makes the state-carrying agent IFR. That result is ungated. If the agent is also gated under the
independent-segment, fresh-retry model of Corollary~\ref{cor:finitefloor} with constant $\sigma$,
Part~I's $1/\sigma$ multiplier is an upper bound for it, up to the rounding above.
Bounded-load holes (R-E case (b)) bound the hazard but do not order it.

\begin{example}[Constant expected arrival rate does not determine a horizon]\label{ex:latent}
Draw $Z$ once: $Z=0$ with probability $0.9$, and $Z=10\lambda$ otherwise, with $\lambda>0$.
Conditional on $Z$, arrivals are homogeneous Poisson of rate $Z$. Then $\E N(T)=\lambda T$,
but $\Prb(N(T)=0)=0.9+0.1e^{-10\lambda T}>0.9$. The half-success horizon is infinite.
An ordinary rate-$\lambda$ Poisson process has the same expected counts and finite horizon
$\log2/\lambda$. For any nonnegative random rate $Z$,
$\E e^{-ZT}\ge e^{-\E[Z]T}$ by Jensen's inequality. This compares zero arrivals at fixed
mean rate, not multi-fault damage or retry cost.
\end{example}

A latent rate changes the horizon value. If visibility does not depend on the latent rate, the
rounded $1/\sigma$ ceiling survives at every segment length, and in the ideal vanishing-segment
limit the multiplier is exactly $1/\sigma$. At finite segment length the exact multiplier is not
invariant: with ideal gates at a deterministic rate $z>0$, log-linear interpolation gives
$H_{\mathrm{gated}}/H_{\mathrm{raw}}=zs/\log(1-\sigma+\sigma e^{zs})$, which depends on $z$ when
$0<\sigma<1$. A numerical ratio also requires finite horizons.

\begin{proposition}[Rate mixing preserves the multiplier]\label{prop:mixing}
Draw $Z\ge0$ once per run. Conditional on $Z$, let the independent-segment model of
Corollary~\ref{cor:finitefloor} hold with constant intensity $Z\lambda_0$ and a stealth share
$\sigma\in(0,1]$ that does not depend on $Z$. Let $S_{\mathrm{raw}}(t)=\E e^{-Z\lambda_0t}$
be the ungated survival, and let segments have equal length $s$. Then, at every segment
boundary $T$,
\[
S_{\mathrm{gated}}(T)\le S_{\mathrm{raw}}(\sigma T),
\qquad\text{hence}\qquad
H_{\mathrm{gated}}\le s\left\lceil H_{\mathrm{raw}}/(\sigma s)\right\rceil
\]
for the discrete horizon measured in work.
For ideal gates ($\alpha=\sigma$, $\beta=0$, as in Proposition~\ref{prop:hazard}), the first
inequality becomes an equality in the vanishing-segment limit. Without ideal gates it need not:
with $\sigma=\alpha=1/2$ and $\beta=1/2$, each accepted-segment survival is exactly
$e^{-Z\lambda_0s}$, so $S_{\mathrm{gated}}(T)=S_{\mathrm{raw}}(T)$. If instead $\sigma=\sigma(Z)$, then
$S_{\mathrm{gated}}(T)\le\E e^{-\sigma(Z)Z\lambda_0T}$, and no fixed multiplier follows.
\end{proposition}
\begin{proof}
Corollary~\ref{cor:finitefloor}, applied conditional on $Z$, gives
$S_{\mathrm{gated}}(T\mid Z)\le e^{-\sigma Z\lambda_0T}$. Take expectations. At the first
boundary $T\ge H_{\mathrm{raw}}/\sigma$, $S_{\mathrm{raw}}(\sigma T)\le S_{\mathrm{raw}}(H_{\mathrm{raw}})=1/2$
by continuity and monotonicity of $S_{\mathrm{raw}}$. For the limit, apply
Proposition~\ref{prop:hazard} conditional on $Z$ and use bounded convergence. An interpolated
bound needs log-linear interpolation of the conditional boundary survivals before averaging.
\end{proof}
Without rounding the horizon bound is false at fixed $s$: the example in the proof of
Proposition~\ref{prop:monotone} is the case $Z\equiv1$.
The survival form also covers Counterexample~\ref{ex:latent}, where both horizons are infinite.
With $Z\in\{0.5,1.5\}$ equally likely, $\lambda_0=1$, and ideal gates, independent visibility
$\sigma=0.2$ gives multiplier $5.0000$ in the limit. Visibility $\sigma(Z)=(0.1,0.3)$, with the
same mean $0.2$, gives $5.0549$. Together with Proposition~\ref{prop:monotone}, for the ideal
vanishing-segment survival curve: random rate mixing keeps the multiplier, deterministic IFR
lowers it, DFR raises it, and visibility that co-varies with difficulty removes any fixed value.
At finite segment length the gated horizon can be shorter: with constant rate $1$ (both IFR and
DFR), ideal gates, $\sigma=1/2$, and $s=1.2$, one segment survives with probability
$2/(1+e^{1.2})<1/2$, so the discrete horizon is $1.2<2\log2$. The DFR lower bound concerns
$H_\sigma$, not finite-segment horizons.

\section{Checkpoint cost and persistent state}\label{sec:cost}
\begin{proposition}[Heterogeneity changes mean retry cost]\label{prop:cost}
Consider equal-length segments $s>0$, perfect detection, no false rejection, independent fresh
retries within each segment, and a segment-specific rate $L\ge0$ that stays fixed during its
retries. Each attempt costs $s+G$. If the workload distribution of rates has finite mean, then
\[
C(s)=(1+G/s)\E[e^{Ls}]\ge(1+G/s)e^{\E[L]s}.
\]
Equality holds only for a constant rate (when $s>0$); the left side can be infinite.
If the rate distribution does not depend on $s$, and the moment and its derivative are finite
near an interior optimizer, the stationarity condition is
\begin{equation}\label{eq:checkpoint}
\frac{\E[L e^{Ls}]}{\E[e^{Ls}]}=\frac{G}{s(s+G)}.
\end{equation}
\end{proposition}
\begin{proof}
A rate-$L$ attempt passes with probability $e^{-Ls}$. Expected attempts are $e^{Ls}$.
Average the conditional cost and apply strict convexity of the exponential. Differentiating
$\log C(s)$ gives the stationarity equation under the stated moment conditions.
\end{proof}
For constant $L=\lambda$, the exact stationary equation in this model is
$s(s+G)=G/\lambda$. Its positive root approaches $\sqrt{G/\lambda}$ when $G\lambda\to0$.
This is the model used for Part I's approximation. Other detection times and recovery costs
give other checkpoint formulas; general failure distributions already have a large
checkpointing literature~\cite{toueg}.

At $s=1,G=0.1$, equal workload shares with rates $0$ and $2$ give
$C=1.1(1+e^2)/2\approx4.6140$. Replacing them by their mean rate $1$ gives
$1.1e\approx2.9901$. The inequality of Proposition~\ref{prop:cost} holds at every $s$, so it
also holds between the two optimal costs. With the two rates held fixed as $s$ varies, the
optimal gap is much smaller: about $1.8530$ at $s\approx0.239$ against $1.7951$ at
$s\approx0.270$, because the heterogeneous optimum uses shorter segments. The fixed-interval
figure alone overstates the practical effect.
No randomness is needed. Take a deterministic intensity $\lambda(t)$ indexed by work position,
and let a retry repeat the same position, so that it keeps the same rate. A segment's rate is
then the average of $\lambda$ over that segment. At each fixed $s$ these averages have mean
equal to the overall mean rate, so the same Jensen comparison holds at every $s$. Their
distribution changes with $s$, however, so Equation~\eqref{eq:checkpoint} does not apply
directly; the derivative must include that dependence.
Part~I's limitations paragraph states that a clustered or time-varying hazard leaves the mean
checkpointing overhead governed by the same $\lambda$ and inflates only its variance. The
countermodel above assumes that a retry keeps its segment's rate, as a retry of the same hard
step would. The paragraph does not state that assumption, but its own mechanism, runs of bad
segments that trigger consecutive recomputations, is consistent with it. Under that assumption
the statement is false and the mean rises.
This example changes no gate behavior and uses fresh retries within each segment.
If the rate is redrawn independently on each attempt instead, perfect detection gives
expected attempts $1/\E[e^{-Ls}]$. The reset law therefore changes even the direction of a
comparison with mean-rate substitution.

\paragraph{Finite state models.}
A latent-state model can include context poisoning, repair, and reset without pretending that
state is an average rate. Let $X_i$ be a finite state before emitted unit $i$ and let
\[
K_i(x,y)=\Prb(I_i=0,X_{i+1}=y\mid X_i=x)
\]
under a Markov assumption that makes $X_i$ sufficient for this transition. $K_i$ is
substochastic. For initial row distribution $\pi$,
\begin{equation}\label{eq:kernel}
S_n=\pi K_1K_2\cdots K_n\mathbf1.
\end{equation}
This follows by summing probabilities of correct state paths. For retries within one unit,
let $R(x,y)$ be the joint probability of rejection and next state $y$, and let $b(x)$ be
immediate bad acceptance. If the spectral radius of $R$ is below one, the eventual bad
probability vector is $(I-R)^{-1}b$, and expected attempts are $(I-R)^{-1}\mathbf1$.
Both follow by summing the convergent series $\sum_{r\ge0}R^r$.
Perfect reset, retained premises, and learning correspond to different kernels. These formulas
do not establish which kernel is appropriate for an agent.

\paragraph{Propagation.}
One escaped root fault can corrupt many descendants. A model of fresh introduction counts
the root once; a model of faulty outputs counts its consequences. For the metric of
Definition~\ref{def:success}, what happens after the first escaped fault cannot change the
already-failed prefix. It can change damage and repair cost. Excitation by a detected or
otherwise unescaped fault is different: it acts before any escape and changes the conditional
prefix risks. If a root fault occurs with probability $p$ and escapes with probability $\sigma$,
and a detected root causes one further faulty attempt that escapes with probability $\sigma$,
the escape probability is $p\sigma+p(1-\sigma)\sigma$, not $p\sigma$. A self-exciting arrival model can
represent some forms of increased subsequent fault activity~\cite{hawkes}, but does not by
itself identify distinct causes, visibility, or semantic propagation.

\section{Mechanization with heterogeneous risks}\label{sec:psm}
At a correct pre-unit history $h$, let $c_i(h)$ be the probability of routing to a mechanism.
Let $o_i(h)$ be the open-path slip conditional on choosing that path.
On the mechanism branch, let $e_i(h)$ be direct escaped-fault probability and $b_i(h)$ the
probability of rejection into fallback. Let $f_i(h)$ be the fallback slip conditional on
that rejection. Direct escape and rejection are disjoint events.

\begin{proposition}[Routing identity and a restricted closure comparison]\label{prop:closure}
At each such history, the slip probability is
\begin{equation}\label{eq:closure}
q_i(h)=(1-c_i(h))o_i(h)+c_i(h)\{e_i(h)+b_i(h)f_i(h)\}.
\end{equation}
Averaging over histories conditional on $A_{i-1}$ gives the $q_i$ of
Proposition~\ref{prop:survival}.
For a fixed exogenous sequence of independent units, suppose routing changes neither later
input laws nor branch risks. Put $d_i=e_i+b_if_i$. Replacing the open path at unit $i$ by the
mechanism improves its survival factor exactly when $d_i<o_i$.
It strictly improves total finite-run survival if the other factors are positive.
For stationary identical units this recovers Part IV's realized-slip criterion and product law.
\end{proposition}
\begin{proof}
Condition first on routing, then on rejection within the mechanism branch. This gives the
identity. Under the additional comparison assumptions, total survival is the product of
fixed local factors. Replacing $1-o_i$ by $1-d_i$ has the stated effect.
\end{proof}

The comparison assumptions matter. A mechanism with lower immediate risk can move execution
to a state with higher later risk. For example, an open first step can have risk $0.1$ and
leave a safe second step. A closed first step can have risk zero but leave a second-step risk
$0.9$. Total survival falls from $0.9$ to $0.1$. Equation~\eqref{eq:closure} remains true;
the claim that local improvement implies global improvement does not.

\paragraph{Work share and risk share.}
For small risks on a fixed independent workload, $1-S_n=\sum_iq_i+O((\sum_iq_i)^2)$.
For units with small $p_i$, fresh retries, ideal acceptance outside a common stealth set
($\alpha_i=\sigma$, $\beta_i=0$), and closed units with zero escape, the first-order remaining
risk is $\sigma\sum_i(1-c_i)p_i$. A stealth floor alone gives only a lower-bound contribution:
with $\sigma=0.1$ but bad acceptance $0.5$, the open slip is $0.5p_i+O(p_i^2)$.
Define the fault-exposure share $c_{\mathrm{risk}}=\sum_i c_i p_i/\sum_i p_i$ when the denominator
is positive. It equals work share only under additional homogeneity or selection assumptions.
Consider one unit with risk $0.01$ and one with risk $0.09$. Closing either one gives 50\%
work closure, but leaves risk $0.09$ or $0.01$. The original PSM formula is appropriate for
its stationary mixing model; volume share alone does not generalize it to this workload.
These risks describe escapes, not confidence scores.

\paragraph{Costs under changing demand.}
For deterministic exogenous call intensity $r(t)$, expected per-call saving $\Delta(t)$,
upkeep rate $u(t)$, fixed lifetime $\tau$, build cost $B$, and validation cost $V$, linear
accounting gives closure benefit
\begin{equation}\label{eq:benefit}
\int_0^\tau[r(t)\Delta(t)-u(t)]dt-B-V.
\end{equation}
Here per-call saving includes the actual expected retry, fallback, and loss costs on each path.
With random lifetime and demand, use the expectation of the stopped total; products of means
need not agree. Part IV's $(r\Delta-u)\tau>B+V$ is the constant-parameter special case.
Selecting regions separately remains valid only with additive costs and no coupling.

\section{Blind sets, evidence, validation, and drift}\label{sec:evidence}
\begin{proposition}[A structural blind set under changing fault laws]\label{prop:blind}
Let $\Sigma$ be a measurable fault set accepted by every gate in a fixed allowed collection.
For every fault distribution $\mu_t$ and every stack from that collection,
its conditional false-accept mass is at least $\mu_t(\Sigma)$.
For two ideal stacks accepting exactly fault sets $\Sigma_F,\Sigma_N$, the combined mass is
$\mu_t(\Sigma_F\cap\Sigma_N)$. It equals the product of the two masses only under independence
of these events under that same $\mu_t$.
\end{proposition}
\begin{proof}
The combined accepted-fault set contains $\Sigma$. Apply monotonicity of measure.
For the two ideal stacks, conjunction is set intersection; independence is precisely the
condition giving the product of masses.
\end{proof}
Part II supplies the inclusion from its representation and zero-false-rejection assumptions.
Its finite false-rejection extension also needs the stated likelihood-ratio control.
Arrival stationarity is not used. A constant numerical floor does require a lower bound on
$\mu_t(\Sigma)$ in the relevant conditional workloads. A changed representation, predicate,
or allowed evidence class can change $\Sigma$ itself.

\begin{proposition}[Twins under dependent observation]\label{prop:twins}
Fix a target distribution and a mechanism class with fault mass $\pi(M)$.
Suppose two mechanisms produce the same law for the full observation record and auxiliary
randomness used by a promotion rule. Their promotion probabilities are equal.
Let $U=u(\text{record},\text{auxiliary randomness})$ for one fixed measurable rule $u$.
If this upper bound satisfies
$\Prb_M\{\pi(M)\le U\}\ge1-\eta$ for every mechanism in the class, then
\[
\Prb_M\Big\{U\ge\sup_{M'\sim M}\pi(M')\Big\}\ge1-\eta,
\]
where $M'\sim M$ denotes equality of those observation laws.
Pointwise twins have equal observation laws under any common exogenous sampling law,
even with dependent samples, provided the joint auxiliary-randomness law is also common.
\end{proposition}
\begin{proof}
The same measurable rule applied to identically distributed records has identical laws.
The common rule is necessary: an oracle bound $U=\pi(M)$ is valid for every mechanism but is
not a function of the record, and it would separate twins of masses $0$ and $1$.
For each $t$ below the supremum, choose a twin with fault mass greater than $t$ and apply
validity for that twin. Transfer the distribution of $U$ to $M$. The events $\{U\ge t\}$
decrease as $t$ increases to the supremum along a sequence, and their intersection is
$\{U\ge\sup_{M'\sim M}\pi(M')\}$. Continuity of measure from above gives the displayed bound.
Pointwise twins agree on every realized sample, so common joint sampling and randomness give
the required equality in law.
\end{proof}
Common adaptive sampling based only on the shared transcript and common independent random
seeds also preserves pointwise twins, by induction on the transcript. Side information about
mechanism text or identity can break this condition. Neither the result nor Part IV's original
twin result forces the mechanism itself to be wrong. They constrain what the specified evidence
can certify. This extension does not validate the separate generalizing-builder hypothesis.

\begin{proposition}[Validation needs conditional exposure]\label{prop:validation}
Let $D_j$ mean that validation item $j$ detects a fault, with exact evidence on that item.
If every candidate above target fault mass $\theta\in(0,1)$ has
\[
\Prb(D_j\mid D_1^c,\ldots,D_{j-1}^c)\ge\theta
\]
whenever the conditioning event has positive probability, then a zero-detection promotion
rule falsely promotes such a candidate with probability at most $(1-\theta)^n$.
Thus $n\ge\log(1/\eta)/[-\log(1-\theta)]$ is sufficient for error at most $\eta$.
\end{proposition}
\begin{proof}
Apply the chain rule to the event of no detections.
\end{proof}
Independent target-distribution samples satisfy this condition. Repeating one randomly drawn
input $n$ times does not: for a deterministic mechanism of fault mass $\theta'$ the chance
of no detection is $1-\theta'$, for every $n\ge1$.
At $\theta'=0.1,n=20$, the two probabilities are $0.9$ and $0.9^{20}\approx0.1216$.
Part IV's sample-size necessity follows within its original class with masses arbitrarily
close to $\theta$ and independent target samples. Necessity is not asserted for arbitrary
adaptive or targeted validation schemes.

\begin{proposition}[Distribution drift and changing fault sets]\label{prop:drift}
For probability measures $D_t,D_0$, use
$\TV(D_t,D_0)=\sup_A|D_t(A)-D_0(A)|$.
If a mechanism's escaped-fault set $B$ is fixed, then
$|D_t(B)-D_0(B)|\le\TV(D_t,D_0)$.
If the set also changes from $B_0$ to $B_t$, then
\[
|D_t(B_t)-D_0(B_0)|\le\TV(D_t,D_0)+D_0(B_t\mathbin\triangle B_0).
\]
\end{proposition}
\begin{proof}
The first statement is the definition of total variation. For the second, add and subtract
$D_0(B_t)$; the difference between masses of two sets is bounded by their symmetric difference.
\end{proof}
The bound needs information about drift; it does not estimate it. With unrestricted shift,
a mechanism correct on all validation inputs can fail on all future inputs.
A watch/reopen policy therefore needs a measurement model and delay analysis, which remain open.

\section{When the earlier results hold}\label{sec:audit}
The companion assessment (\texttt{fault\_process\_assessment.json}) identifies every formal result by its source label and hash.
The following tables summarize the logical dependencies. ``Preserved'' means preserved under
the listed assumptions, not established empirically. ``Replace'' means the formula changes
outside the original model. No published source is edited by this assessment.

Table~\ref{tab:matrix} classifies the 25 results. \emph{Arrival-free} results use no arrival or
stationarity assumption in their stated conclusion; their masses still depend on the current
fault law. \emph{Conditional} results keep the form of their conclusion when uniform or
conditional bounds hold on the reachable histories. \emph{Formula} results are quantitative
formulas derived under constant or stationary parameters (a constant rate, or stationary
Bernoulli risks); their form changes outside the original model. The classes are applied in
this order of precedence. The counts are 14, 10, and 1.
A bullet marks an axis of Section~\ref{sec:axes} whose change can break the result's stated
conclusion: D difficulty, P dependence, R retry, V visibility, I intervention, W drift. Bullets
concern the conclusion itself, not surrounding applications; for example, Part III's
cost-separation comparison is per unit, while translating an output target into a per-unit target
depends on output dependence. The R bullets on the horizon ceiling and stealth decomposition
concern their exact $H_{\mathrm{raw}}/\sigma$ forms; their first-attempt bounds survive any retry
law. R-A is conditional because its displayed form holds whenever
$q_i\ge p\lambda_{\mathrm{mix}}$ on correct prefixes (Proposition~\ref{prop:survival}). The
zero-escape and reach results are conditional for the same reason: a constant correct-prefix
risk keeps their conclusions, even with dependent outputs. A geometric first-failure index $K$
with $I_i=\ind\{i\ge K\}$ gives dependent outputs and exactly $S_n=(1-f)^n$. The classification
is a review judgment recorded with the assessment; it is not a mechanical proof classification.

\begin{table}[!ht]\centering\small
\caption{Classification of the earlier results and the axes that can break them.}\label{tab:matrix}
\begin{tabular}{lllcccccc}
\toprule Part & Result & Class & D & P & R & V & I & W\\\midrule
I & Per-segment slip & conditional &  &  & $\bullet$ &  &  & \\
I & Threshold and overhead & conditional & $\bullet$ & $\bullet$ & $\bullet$ & $\bullet$ &  & \\
I & Horizon ceiling & conditional & $\bullet$ & $\bullet$ & $\bullet$ & $\bullet$ &  & \\
I & Diversity product & arrival-free &  & $\bullet$ &  & $\bullet$ &  & \\
I & Stealth decomposition & conditional & $\bullet$ & $\bullet$ & $\bullet$ & $\bullet$ &  & \\
I & Optimal interval & formula & $\bullet$ &  & $\bullet$ &  &  & \\
II & GG1 inheritance & arrival-free &  &  &  & $\bullet$ &  & $\bullet$\\
II & GG2 fixed point & arrival-free &  &  &  & $\bullet$ &  & $\bullet$\\
II & Exogenous family & arrival-free &  & $\bullet$ &  & $\bullet$ &  & \\
II & Reality grounding & arrival-free &  &  &  & $\bullet$ &  & $\bullet$\\
III & Series-system identity & conditional & $\bullet$ & $\bullet$ &  &  &  & \\
III & Zero-escape corollary & conditional & $\bullet$ & $\bullet$ &  &  &  & \\
III & Cost separation & arrival-free &  &  &  &  &  & \\
III & Association lemma & arrival-free &  & $\bullet$ &  &  &  & \\
III & Correlated-escape bound & arrival-free &  & $\bullet$ &  &  &  & \\
IV & R-A partial closure & conditional & $\bullet$ & $\bullet$ &  & $\bullet$ & $\bullet$ & $\bullet$\\
IV & Evidence twins & arrival-free &  &  &  &  &  & $\bullet$\\
IV & Same-family evidence & arrival-free &  &  &  & $\bullet$ &  & $\bullet$\\
IV & Exact replay & arrival-free &  &  &  &  &  & $\bullet$\\
IV & Accepted-output copy & arrival-free &  &  &  & $\bullet$ &  & \\
IV & Produced undetected faults & arrival-free &  &  &  & $\bullet$ & $\bullet$ & \\
IV & R-C break-even & conditional &  & $\bullet$ &  &  & $\bullet$ & $\bullet$\\
IV & R-D composition & arrival-free &  & $\bullet$ &  &  &  & \\
IV & Reach and residual & conditional & $\bullet$ & $\bullet$ &  &  & $\bullet$ & \\
IV & R-E bounded load & conditional & $\bullet$ &  &  &  & $\bullet$ & \\
\bottomrule
\end{tabular}
\end{table}

\small
\begin{longtable}{p{0.23\textwidth}p{0.70\textwidth}}
\toprule Result & Assessment \\\midrule\endhead
\multicolumn{2}{l}{\textbf{Part I: gated computation}}\\
Per-segment slip & Preserved with constant conditional retry probabilities. For nonfresh retries,
replace $p\alpha_m/P_{\mathrm{adv}}$ by Proposition~\ref{prop:retry}. Counterexample~\ref{ex:retry}.\\
Threshold and overhead & Preserved under uniform conditional gate accuracy and progress:
Theorem~\ref{thm:threshold}. Mean accuracy alone is insufficient, and unbounded difficulty
($p_i\to1$) removes the uniform acceptance bound. The converse needs persistent badness and
conditional escape, rather than a marginal nonzero rate.\\
Horizon ceiling & Replace a fixed $p\sigma$ by conditional exposure bounds
(Proposition~\ref{prop:floor}). A positive marginal floor alone does not imply a finite horizon.
The printed first-attempt bound $s\log2/(p\sigma)$ survives arbitrary retries; the exact
$H_{\mathrm{raw}}/\sigma$ form needs fresh retries (Corollary~\ref{cor:finitefloor}) or an
equally strong conditional survival bound. The $1/\sigma$ multiplier survives rate mixing (Proposition~\ref{prop:mixing}) and is an upper
bound under IFR (Proposition~\ref{prop:monotone}); under DFR the multiplier can exceed it.\\
Diversity product & Preserved under independence of stealth events under the operative fault
law. Family names or temporal independence do not establish this independence.\\
Stealth decomposition & Its mixture identity survives pointwise under the stated conditional
family model. The appended horizon consequence needs persistent exposure. Under fresh Poisson
retries the exact constant-rate bound follows from Corollary~\ref{cor:finitefloor}; a changed
retry law can break it. Its breadth limit needs $\rho<1$ (errata candidate 4).\\
Optimal interval & The constant-rate approximation survives in its cost model. Heterogeneity,
including a deterministic $\lambda(t)$, raises mean cost at every interval and requires
Equation~\eqref{eq:checkpoint} when the rate distribution does not depend on $s$; persistent state
requires a retry model. The printed minimum drops a $1/(1-\beta_m)$ factor (errata candidate 3).\\
\midrule
\multicolumn{2}{l}{\textbf{Part II: generated gates}}\\
GG1 inheritance & Structural inclusion survives without an arrival model, under fixed
representation and soundness/support assumptions. Mass is evaluated under the new fault law.\\
GG2 fixed point & Lower bound survives for sound stacks. Equality additionally requires
reachable coverage outside the blind set. This is an existing achievability qualification,
independent of arrival stationarity.\\
Exogenous family & Intersection identity survives for the ideal stacks. Product requires
independence under the current law; strict reduction requires positive removed mass, not
set noncontainment (errata candidate 5).\\
Reality grounding & Zero covered-scope escape survives if the oracle matches the current
correctness predicate. Exact residual equality off scope requires further conditions
(errata candidate 6).\\
\midrule
\multicolumn{2}{l}{\textbf{Part III: escape, cost, correlation}}\\
Series-system identity & For independent unequal marginals use $1-\prod_i(1-f_i)$.
For arbitrary dependence use conditional risks or Equation~\eqref{eq:marginal}.\\
Zero-escape corollary & Preserved when the correct-prefix risk is constant, $q_i=f$;
independence across outputs is not needed. For varying sequential risks,
summability governs positive infinite survival; a triangular array has vanishing total escape
iff its conditional-risk sum tends to zero.\\
Cost separation & Preserved if the per-unit power law, geometric gate improvement, fixed
prevalence, good-acceptance rate, and cost assumptions remain. No Poisson law is needed.
The translation from output risk to per-unit target must be revised under dependence.\\
Association lemma & Preserved for same-direction monotone functions of independent primitives;
no constant arrival rate is needed~\cite{epw}.\\
Correlated-escape bound & Preserved under that provenance model at fixed marginals. Arbitrary
dependence can violate it. The bound concerns any escape, not simultaneous damage.\\
\midrule
\multicolumn{2}{l}{\textbf{Part IV: progressive semantic mechanization}}\\
R-A partial closure & Original ceiling needs independent stationary mixing and fixed mechanisms.
Use conditional routing risks for history dependence. Work share alone is insufficient.
Its $(1+O(\lambda s))$ factor is necessary for the closed branch (errata candidate 2).\\
Evidence twins & Preserved under equal observation laws. Independent samples can be weakened
as in Proposition~\ref{prop:twins}; changes that distinguish twins are outside that premise.\\
Same-family evidence & Preserved under the representation restriction; the explicit twin-mass
formula also needs the countable scope and finite-patch closure. Drift changes its target mass.
Evidence not decided within $F$, such as exact correctness evidence, can separate $F$-twins.\\
Exact replay & The disagreement bound is pointwise and survives every fixed target law.
It must be recomputed or bounded under shift.\\
Accepted-output copy & The displayed lower bound $p\lambda_F-\E[\delta]$ survives without fresh
retries: under eventual emission a first-attempt stealth fault is emitted at once. Fresh retries
only give the proof's intermediate denominator step. An exact reference gate removes the floor.
No realized-build floor follows.\\
Produced undetected faults & Preserved with the build/input law, preservation of undetected
outputs, and no discarded builds. Repair that changes undetected outputs can break it. Both
parts depend on visibility: detecting previously invisible build faults removes the floor.\\
R-C break-even & Accounting survives with actual expected costs; use Equation~\eqref{eq:benefit}
for varying demand. Conditional exposure gives a sufficient zero-failure sample count
(Proposition~\ref{prop:validation}); the source's necessary count needs its independent target
sampling and near-threshold class, and targeted validation can need far fewer items. Separable
frontier selection also needs additive uncoupled costs.\\
R-D composition & Preserved under residual completeness and monotone provenance.
Without provenance association, completeness still gives a union bound. Missing interaction
faults can invalidate completeness itself.\\
Reach and residual & Exact geometric horizon and $1/(1-c)$ comparison need the stated fixed
branch parameters and a constant correct-prefix risk $q(c)$; independence across outputs is not
needed. Feedback and selected routing require a new comparison.\\
R-E bounded load & Already uses a varying load-dependent Poisson rate. Its bounds survive
with the stated load/intensity bounds and clock. Bounded load alone establishes neither
monotonicity nor clean reset. With nondecreasing load the ungated agent is IFR; if it is also
gated under the model of Corollary~\ref{cor:finitefloor}, Proposition~\ref{prop:monotone} makes
Part~I's $1/\sigma$ an upper bound on the multiplier, up to rounding.\\
\bottomrule
\end{longtable}
\normalsize

\paragraph{Errata candidates in the published papers.}
The assessment found seven items in Parts I and II. Each carries a severity: a \emph{false claim}
within the stated model, a \emph{strengthening} of a correct bound, a \emph{dropped factor}, an
\emph{overstated equality or strictness}, or an \emph{endpoint} restriction. Only item 1 depends on
the arrival model. These qualifications concern the cited versions; this paper does not amend them.

\begin{enumerate}[leftmargin=*,itemsep=2pt]
\item \emph{Part I, limitations: mean checkpoint overhead (false claim, prose).} The claim that a
clustered or time-varying hazard changes only the variance of checkpointing overhead is false
when a retry keeps its segment's rate, including for a deterministic $\lambda(t)$
(Proposition~\ref{prop:cost} and Section~\ref{sec:cost}). The paragraph does not state that
retry assumption; the countermodel adds it. The effect at the respective optimal intervals can be
small.

\item \emph{Part I, horizon ceiling: the $(1+O(\lambda s))$ factor (strengthening).} The theorem
prints $H\le s\log2/(p\sigma)=(\log2/(\lambda\sigma))(1+O(\lambda s))$, from the floor
$S_n\le(1-p\sigma)^n$. That bound is correct. Retaining the acceptance denominator gives
Corollary~\ref{cor:finitefloor}: under Part I's fresh-retry assumptions and its continuous
horizon convention, $H\le H_{\mathrm{raw}}/\sigma$ holds at every segment length, so the factor can
be removed. The discrete horizon rounds up to whole segments. The stealth decomposition prints
the exact form $H_{\mathrm{gated}}\le H_{\mathrm{raw}}/\lambda_{\mathrm{irr}}$, but its proof
substitutes into the first-order ceiling; Corollary~\ref{cor:finitefloor} supplies the exact form
under fresh retries. Without fresh retries the exact form can fail while the printed first-order
bound holds: with $\lambda=s=1$, $\sigma=1/2$, stealth faults accepted, and a rejected nonstealth
fault followed by a good retry, $H\approx1.8246>H_{\mathrm{raw}}/\sigma\approx1.3863$, while
$s\log2/(p\sigma)\approx2.1931$.
Part~IV's closed paths cannot inherit this. With $\varepsilon_{\mathrm{cl}}=p\lambda_{\mathrm{pc}}$
and $0<\lambda_{\mathrm{pc}}\le1$, the function $t\mapsto-\log(1-t(1-e^{-L}))-tL$ is convex on
$[0,1]$ and vanishes at both ends. Hence $-\log(1-p\lambda_{\mathrm{pc}})\le\lambda_{\mathrm{pc}}L$,
and an all-closed run has exact horizon at least $H_{\mathrm{raw}}/\lambda_{\mathrm{pc}}$. At
$\lambda s=1$ and $\lambda_{\mathrm{pc}}=1/2$ it is $1.3162\,H_{\mathrm{raw}}/\lambda_{\mathrm{pc}}$.
The $(1+O(\lambda s))$ factor of R-A is therefore necessary for its closed component; the open
component can use Corollary~\ref{cor:finitefloor}.

\item \emph{Part I, optimal interval: the minimal overhead (dropped factor).} The printed cost is
$C(s)\approx(1+G/s)e^{\lambda s}/(1-\beta_m)$, but the printed minimum $1+2\sqrt{G\lambda}$ omits
the factor $1/(1-\beta_m)$. For fixed $\beta_m=1/4$ the minimum tends to $4/3$, not $1$, as
$G\lambda\to0$. The corrected statement is
$C_{\min}=(1+2\sqrt{G\lambda}+O(G\lambda))/(1-\beta_m)$, or the theorem should assume small
$\beta_m$. The minimizing interval is unaffected.

\item \emph{Part I, stealth decomposition: convergence (endpoint).} The limit
$\lambda_{\mathrm{irr}}+(1-\lambda_{\mathrm{irr}})\rho^k\to\lambda_{\mathrm{irr}}$ needs $\rho<1$
(or $\lambda_{\mathrm{irr}}=1$); at $\rho=1$ breadth does not help.

\item \emph{Part II, exogenous family: strict reduction (overstated strictness).} Set
noncontainment does not imply a strict reduction in mass. Under uniform measure on $[0,1]$,
$\Sigma_F=[0,1/2]$ and $\Sigma_N=[0,1/2)$ satisfy $\Sigma_F\not\subseteq\Sigma_N$, but both masses
are $1/2$. The condition should be $\mu(\Sigma_F\setminus\Sigma_N)>0$. Part~II's model section states that equality between fault sets may be read up to $\mu$-null
sets; extending that reading to containment also repairs the statement, so this item is a wording
repair.

\item \emph{Part II, grounding proposition: off-scope residual (overstated equality).} The
proposition guarantees that no fault passes on the covered scope. An oracle declared arbitrary off
scope can reject additional stealth faults there, so the displayed off-scope residual mass is not
always an equality. Conversely, a nonideal family stack can pass additional nonstealth faults off
scope. If the family stack accepts exactly $\Sigma_F$ and the oracle accepts everything outside
$S$, residual escape is exactly $\mu(\Sigma_F\setminus S)$. Under the first condition alone it is
at most that mass. The covered-scope guarantee is unaffected.

\item \emph{Part II, GG2: equality needs reachable coverage (overstated equality; already flagged).}
The equality needs sufficient reachable gates to reject faults outside the stealth set. The
published repository already flags this. Its lower-bound proof uses sound gates. Separate
per-gate false-rejection tolerances do not automatically give a fixed tolerance for an
arbitrarily large conjunction.
\end{enumerate}

\paragraph{What has not been proved.}
We do not prove that real faults follow any listed process, that resets remove hidden state,
that real gate stacks satisfy a uniform conditional edge, or that drift is observable in time.
Part IV's generalizing-builder and load-monotonicity hypotheses remain hypotheses.
An arbitrary adaptive provenance graph need not satisfy the monotone primitive model.
The evidence ladder and typed-evidence framing in Part III are not probability theorems and
do not become invalid merely because an arrival model changes.

\section{Executable checks and counterexamples}\label{sec:checks}
The companion script (\texttt{simulate\_fault\_process.py}) enumerates small probability spaces, checks closed forms against direct
enumeration, and evaluates the examples below. Checks are deterministic. They do not establish
the universal propositions; the proofs do that subject to their stated assumptions.
Negative controls deliberately reinstate false independence, pooled-retry, sample-count, or
fixed-multiplier claims and must be rejected by the check runner.

\begin{center}\small
\begin{tabular}{p{0.49\textwidth}rr}
\toprule Comparison & First model & Second model\\\midrule
100 outputs, marginal risk 0.01: independent / shared survival & 0.3660 & 0.9900\\
Persistent retry example: pooled prediction / actual bad fraction & 0.3333 & 0.5000\\
Same example: pooled prediction / actual expected attempts & 1.3333 & 1.5000\\
Mean rate / unequal rates: overhead, $s=1,G=0.1$ & 2.9901 & 4.6140\\
20 validation items, risk 0.1: independent / repeated no-detection & 0.1216 & 0.9000\\
Close half of two units: close low-risk / close high-risk, remaining risk & 0.0900 & 0.0100\\
Optimal overhead, $G=0.1$: mean rate / unequal rates & 1.7951 & 1.8530\\
Rate mixing, $\sigma=0.2$: multiplier, independent / $Z$-dependent visibility & 5.0000 & 5.0549\\
Weibull, $\sigma=0.2$: multiplier, IFR $k=2$ / DFR $k=1/2$ & 2.2361 & 25.0000\\
Closed path, $\lambda s=1$, $\lambda_{\mathrm{pc}}=1/2$: horizon in units of $H_{\mathrm{raw}}/\lambda_{\mathrm{pc}}$, first-order / exact & 1.0000 & 1.3162\\
\bottomrule
\end{tabular}
\end{center}

The assessment manifest binds each prior-result row to the exact source snapshot. A coverage
check rejects a missing formal result or a changed source hash. This is a completeness and
provenance check, not an automated judgment of the source proofs. The local verification
receipt records check output, negative-control exits, and source hashes. No empirical dataset
is introduced.

\section{Related work and limitations}
The constant-hazard baseline is motivated by Ord's task-duration analysis~\cite{ord}.
Hoeffding supplies the exponential majority-error bound~\cite{hoeffding}.
Association and the fixed-marginal reliability bound are classical~\cite{epw}.
Checkpoint placement for nonidentical tasks and general failure distributions predates this
series~\cite{toueg}. Hawkes processes provide one established family for history-dependent
arrivals~\cite{hawkes}. The IFR and DFR classes are standard in reliability theory~\cite{bp}.
Frailty models show that mixing individual hazards changes the population hazard~\cite{vaupel};
Proposition~\ref{prop:mixing} is the corresponding statement for the gating multiplier.
Nonhomogeneous Poisson models of software failures~\cite{goel}, imperfect repair~\cite{brownproschan},
and Markov-modulated Poisson processes~\cite{mmpp} supply the other candidate families of
Table~\ref{tab:catalog}. None is introduced here as a new stochastic process.

Our contribution is an assumption-separated application to the interaction of imperfect gates,
retries, representation floors, and mechanization. The mathematical results are elementary.
The literature check is scoped, not a systematic novelty review. This is a preprint; it has not been
peer reviewed.

The analysis uses an all-correct prefix metric, eventual emission except in explicit
counterexamples, and declared clocks. It does not infer an arrival process from benchmark
success rates. It does not optimize a general adaptive policy, price correlated catastrophic
loss, or analyze a watch policy's detection delay. Numerical parameters are illustrative.
Changing a model outside its assumptions does not overturn a result within those assumptions.

\section{Conclusion}
Constant fault intensity is not the common foundation of every result in the series.
Of the 25 named results, 14 are arrival-free, 10 survive under stated conditional bounds, and 1
is a stationary-parameter formula that needs replacement outside its original model.
Conditional accuracy and progress suffice for a logarithmic verification threshold.
For the ideal gated survival curve, the $1/\sigma$ ceiling multiplier survives random rate
mixing, is an upper bound under an increasing failure rate, and can be exceeded under a decreasing
one; at finite segment length the upper bounds hold after rounding.
Persistent conditional exposure, rather than a marginal floor alone, creates a horizon ceiling.
Blind-set and evidence-indistinguishability results survive many changes in arrivals, while
their numerical masses depend on the workload. Mechanization changes both exposure and
dependence. Its effect requires branch risks, fallback, state transitions, and maintenance
costs to be stated. These distinctions specify when Parts I--IV apply and where a different
model is required.

\section*{Code and data availability}
\begin{sloppypar}
The standard-library Python script \texttt{simulate\_fault\_process.py} and the assessment
\texttt{fault\_process\_assessment.json} reproduce the checks and tables of Section~\ref{sec:checks} and the
result classification of Section~\ref{sec:audit}. Both are at \url{https://github.com/ivmat/gated-computation-sim}.
This paper is archived on Zenodo, DOI~\href{https://doi.org/10.5281/zenodo.23293740}{10.5281/zenodo.23293740}.
\end{sloppypar}

\begin{thebibliography}{99}
\bibitem{p1} I. Matija\v{s}evi\'c. A Fault-Tolerance Threshold for Gated Agentic Computation.
2026. \url{https://doi.org/10.5281/zenodo.20820968}.
\bibitem{p2} I. Matija\v{s}evi\'c. Generated Gates Inherit Their Generator's Blind Spots.
2026. \url{https://doi.org/10.5281/zenodo.20837102}.
\bibitem{p3} I. Matija\v{s}evi\'c. Escape, Cost, and Correlation at the Verification Floor of
Gated Agentic Computation. 2026. \url{https://doi.org/10.5281/zenodo.21456783}.
\bibitem{p4} I. Matija\v{s}evi\'c. Progressive Semantic Mechanization. 2026.
\url{https://doi.org/10.5281/zenodo.23119225}.
\bibitem{ord} T. Ord. Is there a half-life for the success rates of AI agents? 2025.
\url{https://arxiv.org/abs/2505.05115}.
\bibitem{hoeffding} W. Hoeffding. Probability inequalities for sums of bounded random variables.
\emph{Journal of the American Statistical Association} 58(301):13--30, 1963.
\url{https://doi.org/10.1080/01621459.1963.10500830}.
\bibitem{epw} J. D. Esary, F. Proschan, and D. W. Walkup. Association of random variables,
with applications. \emph{Annals of Mathematical Statistics} 38(5):1466--1474, 1967.
\url{https://doi.org/10.1214/aoms/1177698701}.
\bibitem{toueg} S. Toueg and \"O. Babaoglu. On the optimum checkpoint selection problem.
\emph{SIAM Journal on Computing} 13(3):630--649, 1984.
\url{https://doi.org/10.1137/0213039}.
\bibitem{hawkes} A. G. Hawkes. Spectra of some self-exciting and mutually exciting point processes.
\emph{Biometrika} 58(1):83--90, 1971.
\url{https://doi.org/10.1093/biomet/58.1.83}.
\bibitem{bp} R. E. Barlow and F. Proschan. \emph{Mathematical Theory of Reliability}.
Wiley, New York, 1965.
\bibitem{vaupel} J. W. Vaupel, K. G. Manton, and E. Stallard. The impact of heterogeneity in
individual frailty on the dynamics of mortality. \emph{Demography} 16(3):439--454, 1979.
\url{https://doi.org/10.2307/2061224}.
\bibitem{goel} A. L. Goel and K. Okumoto. Time-dependent error-detection rate model for software
reliability and other performance measures. \emph{IEEE Transactions on Reliability}
R-28(3):206--211, 1979.
\bibitem{brownproschan} M. Brown and F. Proschan. Imperfect repair.
\emph{Journal of Applied Probability} 20(4):851--859, 1983.
\bibitem{mmpp} W. Fischer and K. Meier-Hellstern. The Markov-modulated Poisson process (MMPP)
cookbook. \emph{Performance Evaluation} 18(2):149--171, 1993.
\end{thebibliography}
\end{document}
